\chapter{Construcción del conjunto de datos y del modelo}
\label{cap:dataset-modelo}

\section{Selección del vocabulario}

\pendiente{Justificar las señas elegidas y su cantidad. Validar contra el vocabulario vigente al cerrar el entrenamiento.}

\section{Protocolo de captura}

\pendiente{Describir cómo se capturaron las muestras: personas, repeticiones, condiciones de luz, posición de cámara, herramienta de captura y cantidad total de muestras.}

\section{Representación de las señas}

\pendiente{Explicar la representación de cada seña como secuencia de $T$ cuadros con 126 valores por cuadro (21 puntos por mano, dos manos, tres coordenadas), es decir, una matriz $\mathbf{X} \in \mathbb{R}^{T \times 126}$.}

\section{Preprocesamiento y normalización}

\pendiente{Describir la normalización temporal a longitud fija, el tratamiento de manos ausentes y cualquier normalización espacial aplicada.}

\section{Arquitectura del modelo}

\pendiente{Describir la arquitectura final y justificarla. Contrastar con la arquitectura LSTM planteada en la propuesta si hubo cambios.}

\section{Entrenamiento}

\pendiente{Partición de datos, hiperparámetros, regularización y curvas de entrenamiento.}

\section{Conversión a TensorFlow Lite y cuantización}

\pendiente{Describir la conversión del modelo y comparar variantes (por ejemplo, float32 frente a optimizada) en tamaño y desempeño.}
